\documentclass[11pt]{article}
\usepackage[a4paper,margin=25mm]{geometry}
\usepackage{amsmath,amssymb,amsthm,booktabs,array,longtable}
\usepackage[T1]{fontenc}
\usepackage{lmodern,microtype}
\usepackage[colorlinks=true,linkcolor=blue,citecolor=blue,urlcolor=blue]{hyperref}
\usepackage{enumitem}
\setlist{nosep}
\setlength{\parskip}{5pt}
\setlength{\parindent}{0pt}
\setlength{\emergencystretch}{2em}
\newtheorem{proposition}{Proposition}
\newtheorem{lemma}{Lemma}
\newtheorem{assumption}{Assumption}
\newcommand{\norm}[1]{\left\lVert #1\right\rVert}
\newcommand{\ip}[2]{\left\langle #1,#2\right\rangle}
\newcommand{\T}{\mathbb T}
\newcommand{\dd}{\mathrm d}
\title{State-dependent splitting errors for anisotropic transport and diffusion:\\Counterexamples and a conditional corrected bound}
\author{Independent technical note\\
\small Prepared with OpenAI Codex; published by the GitHub account Teylersf\\
\small AI-assisted draft, not peer reviewed}
\date{28 October 2026}
\begin{document}
\maketitle
\begin{abstract}
We audit the nine-page arXiv version 2603.08799v1 of Zylberman, Fredon,
Loureiro, and Debbasch. Two explicit smooth periodic examples contradict,
respectively, the assertion that all derivative norms are conserved by the
allowed transport flows and the printed leading diffusion splitting
coefficient. The diffusion example has strictly positive coefficients and
positive, nonzero-mean initial data. We derive the missing mixed derivatives,
explain the distinction between a loose operator-norm comparator and an
algorithmic speedup, and document regularity, normalization, and resource
assumptions that need further justification. As a replacement framework,
we prove an exact autonomous finite-dimensional splitting bound with an
explicit uniform commutator hypothesis and a normalized-state corollary.
Deterministic classical checks accompany the note. We neither claim to have
found every error nor prove that mesh-independent vector-norm scaling is
impossible. This is a correction and reproducibility study, not a new
end-to-end quantum algorithm or a claim of priority.
\end{abstract}

\section{Scope, provenance, and standards of evidence}
The audited object is the local nine-page PDF identifying itself as
arXiv:2603.08799v1, submitted 9 March 2026 \cite{original}. Its SHA-256 is
\texttt{d03a710413b32ac622700adcfbaa9c3b860}\\
\texttt{4ab9db00e45c348ea06c3d4be8988}.
The source record links a published conference chapter \cite{chapter};
we checked its public metadata, but did not compare the full chapter text.
No finding here is asserted to apply unchanged to an unexamined revision,
a later journal treatment, or other work by the same authors.

We distinguish four evidence categories:
\begin{enumerate}
\item \textbf{Demonstrated error:} an identity, inequality, or prescription
fails as written, with an explicit calculation.
\item \textbf{Missing hypothesis or proof gap:} the intended result may hold
under additional assumptions, but the displayed argument does not establish it.
\item \textbf{Discretization limitation:} a property of the selected scheme
that does not necessarily destroy its intended convergence regime.
\item \textbf{Editorial ambiguity:} notation or indexing that should be repaired
without implying a substantive mathematical failure.
\end{enumerate}
The motivation for vector-norm analysis remains sound. An, Fang, and Lin
previously studied circumstances in which state-dependent Trotter estimates
avoid large discretized operator norms \cite{an}. A failed prefactor does not
by itself invalidate that research program.

\section{Model and consistent notation}
On the unit periodic torus $\T^d$, consider separately
\begin{align}
\partial_t f+\sum_{j=1}^d c_j(x,t)\partial_j f&=0,
\label{eq:transport}\\
\partial_t u=\sum_{j=1}^d \kappa_j(x,t)\partial_j^2u .
\label{eq:diffusion}
\end{align}
The source imposes $\partial_jc_j=\partial_j\kappa_j=0$.
Under this restriction, \eqref{eq:diffusion} equals the divergence-form
equation in its Eq.~(1). For forward diffusion we explicitly require
$\kappa_j\geq0$. The stronger ellipticity condition
$\kappa_j\geq\kappa_*>0$ is useful for some regularity arguments, but is
not needed for the elementary finite-grid contraction below.

Let $M_j=2^{n_j}$, $h_j=M_j^{-1}$, and $M=\prod_jM_j$.
The total number of data qubits is $n_{\rm tot}=\sum_j n_j$.
A grid coordinate is $x_j=r_j/M_j$, $0\leq r_j<M_j$;
$r_j$ is the integer register label. We use the quadrature norm
\[
 \norm{v}_h^2=M^{-1}\sum_{\boldsymbol r}|v_{\boldsymbol r}|^2.
\]
This converges to the continuum $L^2$ norm for smooth data on the unit torus.
The associated unit quantum state has amplitudes
$v_{\boldsymbol r}/(\sqrt M\,\norm{v}_h)$.
An unnormalized grid solution and its normalized quantum encoding are
different objects and will not share a symbol silently.

For centered first differences write
\[
 \delta_{p,h}v=h^{-1}\sum_{k=-p}^p a_k S^k v,\qquad
 a_{-k}=-a_k,\qquad D_{p,h}=-i\delta_{p,h}.
\]
Periodic shifts make $D_{p,h}$ Hermitian. Because the coefficient in
direction $j$ is independent of that direction, its multiplication matrix
commutes with $D_j$. Consequently $-ic_jD_j$ is skew-Hermitian and
$-\kappa_jD_j^2$ is Hermitian negative semidefinite when $\kappa_j\geq0$.
The source's directional exponentials then yield unitary transport
factors and contractive, generally non-unitary diffusion factors.

\section{A transport counterexample and energy estimate}
\subsection{Derivative norms are not generally conserved}
Page 6, footnote 10, asserts conservation of the $L^2$ norm of any
spatial derivative of $f$. The coefficient conditions stated on page 2
do not imply this. Choose the allowed time-independent periodic shear
\[
 c(x,y)=(\sin(2\pi y),0),\qquad f_0(x,y)=\sin(2\pi x).
\]
Both components are smooth and independent of their own coordinate.
The exact solution is
\[
 f(x,y,t)=\sin\bigl(2\pi[x-t\sin(2\pi y)]\bigr).
\]
Direct differentiation and integration give
\begin{align}
 \norm{f(t)}_{L^2}&=1/\sqrt2,\\
 \partial_y f&=-4\pi^2t\cos(2\pi y)
       \cos\bigl(2\pi[x-t\sin(2\pi y)]\bigr),\\
 \norm{\partial_yf(t)}_{L^2}&=2\pi^2|t|.
\end{align}
The derivative norm starts at zero and grows. This is an analytic
counterexample, not a conclusion inferred from a numerical plot.

\subsection{What can replace conservation}
The velocity is divergence-free because $\sum_j\partial_jc_j=0$.
For a smooth solution, differentiation of \eqref{eq:transport} gives
\[
 (\partial_t+c\cdot\nabla)\nabla f=-(\nabla c)^{\mathsf T}\nabla f,
\]
where $(\nabla c)_{ij}=\partial_jc_i$. Periodic integration
by parts yields
\[
 \frac{\dd}{\dd t}\norm{\nabla f(t)}_{L^2}^2
 \leq 2\norm{\nabla c(t)}_{L^\infty,\mathrm{op}}
                 \norm{\nabla f(t)}_{L^2}^2 .
\]
Thus
\[
 \norm{\nabla f(t)}_{L^2}\leq
 \exp\!\left(\int_0^t\norm{\nabla c(s)}_{L^\infty,\mathrm{op}}\dd s\right)
 \norm{\nabla f_0}_{L^2}.
\]
This bound involves the initial full gradient and the velocity gradients;
it does not conserve each coordinate derivative separately.
Higher Sobolev estimates require corresponding coefficient derivatives.
Constants can therefore depend on the flow and on $T$, even when they
are independent of $h$.

\subsection{Directional commutators}
For $L_j=c_j\partial_j$ and $L_m=c_m\partial_m$,
\[
 [L_j,L_m]f=
 c_j(\partial_j c_m)\partial_mf-
 c_m(\partial_mc_j)\partial_jf.
\]
The transport prefactor on page 6 has this appropriate first-derivative
structure. Its use of a supremum over evolved derivatives can remain
meaningful. The invalid step is identifying those derivative norms with
their initial coordinatewise values.

\section{Diffusion: missing terms and a counterexample}
\subsection{Full continuum commutator}
Put $A_j=\kappa_j\partial_j^2$. Expanding products of differential operators
for $j\ne m$ gives
\begin{align}
[A_j,A_m]u={}&
 \kappa_j(\partial_j^2\kappa_m)\partial_m^2u
 +2\kappa_j(\partial_j\kappa_m)\partial_j\partial_m^2u \notag\\
 &-\kappa_m(\partial_m^2\kappa_j)\partial_j^2u
 -2\kappa_m(\partial_m\kappa_j)\partial_m\partial_j^2u .
 \label{eq:comm}
\end{align}
The common fourth-derivative term cancels. The mixed third derivatives
do not cancel under $\partial_j\kappa_j=0$.
The displayed $a'_g$ following the source's Eq.~(8), page 6, contains only
the first and third terms of \eqref{eq:comm} in its norm estimate.

A valid direct continuum estimate is
\begin{align}
\norm{[A_j,A_m]u}_{L^2}\leq{}&
 \norm{\kappa_j}_\infty\norm{\partial_j^2\kappa_m}_\infty
       \norm{\partial_m^2u}_{L^2}\notag\\
 &+2\norm{\kappa_j}_\infty\norm{\partial_j\kappa_m}_\infty
       \norm{\partial_j\partial_m^2u}_{L^2}\notag\\
 &+\norm{\kappa_m}_\infty\norm{\partial_m^2\kappa_j}_\infty
       \norm{\partial_j^2u}_{L^2}\notag\\
 &+2\norm{\kappa_m}_\infty\norm{\partial_m\kappa_j}_\infty
       \norm{\partial_m\partial_j^2u}_{L^2}.
 \label{eq:corrected-comm}
\end{align}
It requires control of additional derivatives. This is an identity and
triangle-inequality bound in the continuum, not by itself a theorem about
all discrete split trajectories.

\subsection{Strictly positive coefficients and positive initial data}
On $\T^2$ choose, for $0<a,b<1$ and integer $m\geq2$,
\[
 \kappa_x=1,\qquad \kappa_y=1+a\sin(2\pi x),\qquad
 u_0=1+b\cos(2\pi m x)\cos(2\pi y).
\]
The coefficients satisfy all source coordinate restrictions and are
strictly positive. The initial data are positive with mean 1.
Let $q=2\pi$, $A=\partial_x^2$, $B=\kappa_y\partial_y^2$.
Then, with the commutator convention $[A,B]=AB-BA$,
\[
 [A,B]u_0=abq^4
 \left[(m+\tfrac12)\sin(q(m+1)x)
 +(m-\tfrac12)\sin(q(m-1)x)\right]\cos(qy).
\]
Fourier orthogonality gives
\[
 \norm{u_0}_{L^2}^2=1+b^2/4,\qquad
 \norm{[A,B]u_0}_{L^2}
 =\frac{abq^4}{2\sqrt2}\sqrt{4m^2+1}.
\]
The quantity built from the two retained terms in the source, before
the common factor $1/2$ and normalization, is only
\[
 \norm{\kappa_x}_\infty
 \norm{\partial_x^2\kappa_y}_\infty
 \norm{\partial_y^2u_0}_{L^2}
 =abq^4/2 .
\]
The ratio of the actual commutator norm to that printed comparator is
\[
 R_m=\sqrt{(4m^2+1)/2}>1 .
\]
For $m=4$, $R_m=\sqrt{65/2}\simeq5.700877$.

\subsection{Why normalization does not remove this counterexample}
For real self-adjoint $A,B$, the commutator is skew-adjoint, so
$\ip{u_0}{[A,B]u_0}=0$. The differential of $N(v)=v/\norm v$ at $u_0$
is
\[
 DN_{u_0}(w)=
 \frac{(I-\hat u_0\hat u_0^{\mathsf T})w}{\norm{u_0}},
 \qquad \hat u_0=u_0/\norm{u_0}.
\]
The leading splitting defect lies perpendicular to $u_0$, and survives
normalization unchanged apart from division by $\norm{u_0}$.

To connect this continuum example to the source's actual stencil, set
$A_\eta=-D_x^2$ and $B_\eta=-K_yD_y^2$ on sufficiently fine centered
periodic grids. Acting on this fixed trigonometric data and coefficient
family, the discrete derivatives and their commutators converge in
quadrature norm to the displayed continuum expressions. The printed
initial-time comparator converges to the smaller expression above.
Both matrices are real self-adjoint, so the discrete orthogonality also
holds exactly. Thus the strict ratio persists on sufficiently fine
fixed grids, before taking the time-step limit.

On each such fixed finite grid, the global first-order Lie-splitting expansion
is an ordinary matrix Taylor calculation:
\[
 \lim_{L\to\infty}L(S_{T/L}^{\,L}-U(T))u_0
 =\frac T2\int_0^T U(T-s)[B,A]U(s)u_0\,\dd s ,
 \label{eq:global-asymp}
\]
where $U(t)=e^{t(A+B)}$ and $S_h=e^{hB}e^{hA}$.
Here $A,B$ denote the fixed-grid matrices just defined. Their autonomous
exponentials are analytic, so the matrix expansion is justified.
As $T\downarrow0$, the normalized leading coefficient divided by $T^2$
approaches $\norm{[B,A]u_0}/(2\norm{u_0})$.
The printed derivative-only prefactor approaches the smaller comparator
above divided by $2\norm{u_0}$.
Thus the claimed leading coefficient fails for sufficiently small
positive $T$, if the remainder is genuinely $o(1/L)$ as the leading-error
interpretation requires. If a remainder is allowed to contain another
$1/L$ contribution, the advertised prefactor has not been isolated.
This argument concerns a coefficient, not an impossibility of
mesh-independent bounds with \eqref{eq:corrected-comm} and adequate regularity.

\section{Other errors, limitations, and unresolved proof obligations}
\subsection{Positivity is an essential explicit hypothesis}
Page 5, footnote 8, uses contraction for diffusion. This follows from
$\kappa_j\geq0$, but that inequality is not among the explicit mathematical
assumptions on page 2. Negative constant $\kappa=-1$ satisfies the written
smoothness and coordinate restrictions and gives backward heat:
$u=e^{(2\pi)^2t}\cos(2\pi x)$.
This is a missing hypothesis; it is not a criticism of physically
nonnegative thermal conductivity.

\subsection{Squared centered derivatives have extra zero modes}
For $p=1$,
\[
 -\delta_{1,h}^2u_r
 =-\frac{u_{r+2}-2u_r+u_{r-2}}{4h^2},
 \qquad \lambda(\theta)=\sin^2\theta/h^2 .
\]
Besides the constant mode, the checkerboard $u_r=(-1)^r$ lies in the
kernel on an even grid. For every antisymmetric centered first-derivative
stencil, $\sum_k a_k e^{ik\pi}=0$.
In $d$ dimensions the tensor parity modes supply at least $2^d$ null
modes, including with the source's admissible variable coefficients.
This undermines uniform damping and some interpretations of long-time
equilibrium. It does not contradict convergence on a fixed sufficiently
smooth resolved input. A standard nearest-neighbor negative second
difference has symbol $4\sin^2(\theta/2)/h^2$ and removes this extra
one-dimensional zero. Its positivity and quantum implementation would
need a separate analysis; substitution is not a free algorithmic repair.

\subsection{Spatial residuals need stronger derivative control}
The source's page-5 diffusion discussion describes a constant in terms
of derivatives through order $2p+1$. But squaring the first derivative
gives, for $p=1$,
\[
 \delta_{1,h}^2u=u''+\frac{h^2}{3}u^{(4)}+O(h^4).
\]
The direct $O(h^{2p})$ pointwise consistency argument for the squared
stencil needs order $2p+2$ control. Parabolic smoothing or an integrated
estimate might weaken initial regularity requirements, but such an
argument is not supplied. We classify this as a regularity/proof gap,
rather than exclude every possible low-regularity convergence theorem.

\subsection{The directional error budget must sum to the target}
The page-5 bound has the form
$TK\sum_j\norm{c_j}_\infty h_j^{2p}$.
The printed prescription for each $n_j$ assigns an error budget
$\epsilon$ to every term. It ensures at most $d\epsilon$ from that
inequality, not $\epsilon$. Choose $\epsilon_j>0$ with
$\sum_j\epsilon_j\leq\epsilon$, and use
\[
 n_j\geq
 \max\!\left\{0,\left\lceil\frac1{2p}
 \log_2\frac{TK\norm{c_j}_\infty}{\epsilon_j}\right\rceil\right\}
\]
for nonzero coefficients, with zero-coefficient directions handled
separately and any practical minimum grid imposed additionally.
An equal split uses $\epsilon_j=\epsilon/d$.

\subsection{Normalization and sampling cannot be conflated}
Continuum transport conserves its integral $L^2$ norm. Its values sampled
on a fixed finite grid need not have exactly constant quadrature norm
as the continuum profile evolves. A normalized sampled state obeys an
extra derivative-of-normalizer term. Alternatively one can derive a
residual equation for fixed-normalizer samples and normalize only after
comparison. For diffusion, writing the unnormalized ODE solution as a
normalized qubit state introduces a similar distinction. Factors involving
the final norm must be retained in a normalized error estimate.

A useful algebraic repair, valid in any Hilbert norm, is
\begin{lemma}[Normalization stability]
For nonzero $v,w$,
\[
 \left\|\frac v{\norm v}-\frac w{\norm w}\right\|
 \leq\frac{2\norm{v-w}}{\norm w}.
\]
\end{lemma}
\begin{proof}
Insert $v/\norm w$, use the triangle inequality, and apply
$|\norm v-\norm w|\leq\norm{v-w}$.
\end{proof}
For mean-preserving diffusion on the unit torus,
$\norm{u(T)}_{L^2}\geq|\int_{\T^d}u_0|$.
For a quantum-normalized input the corresponding bound must divide
this mean by $\norm{u_0}_{L^2}$.
The source does impose a nonzero-mean condition on page 7; its later
appearance is a scope restriction, not an undiscovered omission.
No positive uniform lower bound follows from this route for mean-zero
input.

\subsection{Remainders and time dependence}
A fixed-grid expansion with an $O_h(\Delta t^3)$ remainder is insufficient
to establish a bound uniform over grids. Its constant may grow as
$h\to0$. Mixed third derivatives in \eqref{eq:comm}, higher nested
commutators, stability of partial split trajectories, and time derivatives
of coefficients all need suitable uniform control.
An $O(T^2/L)$ expression also does not isolate polynomial dependence
on $T$ if its prefactor itself contains evolved Sobolev norms growing
with $T$.

\subsection{What the exponential comparator does and does not show}
Writing $M_{\rm axis}=2^n$ on an equal grid, the source's crude product
bound uses $O(M_{\rm axis}^2)$ for transport and
$O(M_{\rm axis}^4)$ for diffusion commutators.
It is not a lower bound on the work actually needed by an algorithm.
For fixed smooth coefficients and a fixed stencil, cancellation already
improves a direct operator estimate. Indeed
\[
 [C_jD_j,C_mD_m]
 =C_j[D_j,C_m]D_m-C_m[D_m,C_j]D_j,
\]
and a Lipschitz coefficient gives $\norm{[D_j,C_m]}=O(1)$.
Hence this operator-norm commutator bound is $O(M_{\rm axis})$.
For diffusion,
$\norm{[D_j^2,K_m]}=O(M_{\rm axis})$ and
\[
[K_jD_j^2,K_mD_m^2]
=K_j[D_j^2,K_m]D_m^2-K_m[D_m^2,K_j]D_j^2.
\]
The corresponding bound is
$O(M_{\rm axis}^3)$.
These are upper bounds, not asserted sharp rates in every case.
Constant coefficients commute, giving no directional splitting error
in the autonomous case.
Therefore the advertised $\Theta(4^n)$ and $\Theta(16^n)$ improvements
must be read as comparisons with a particular loose bound, not
universal tight speedup factors over all operator-norm analyses.

\subsection{Quantum resource and readout assumptions}
A non-unitary contraction is not a deterministic ordinary unitary gate
on a closed data register. An implementation must specify a dilation,
block encoding, measurement/postselection, amplitude amplification, or
another realization with its associated costs.
For an ideal contraction block with normalization factor 1, success
probability on a normalized input is $\norm{Vu}^2$.
Other normalizations or implementations change this cost.

Merely being differentiable does not uniformly bound the size of a
Fourier or Walsh approximation, its coefficient-access cost, or the
cost of loading arbitrary initial data. The finite-difference symbol
has amplitude growing with grid refinement. At fixed stencil order its
oscillation count as a function of normalized frequency does not grow,
but the associated unitary exponential can become increasingly
oscillatory at a fixed time. A gate-level complexity theorem must
quantify those dependencies rather than appeal to differentiability
alone.

For a fixed observable, an expectation $\ip{\psi}{O\psi}$ is invariant
under a global phase. In particular $f$ and $-f$ encode the same quantum
ray, whereas their signed point values differ. A position projector
returns a normalized squared magnitude, not the signed physical
field value. Recovering amplitudes requires an appropriate reference,
controlled state preparation or overlap protocol, a phase convention,
normalization information, and measurement-cost accounting.
The measurement discussion on page 5 is therefore incomplete when
interpreted as direct field-value recovery by a single expectation.

\section{An exact conditional replacement theorem}
We deliberately state a finite-dimensional autonomous result. It does
not silently extend to the source's time-dependent PDEs.

For each grid $\eta$, let $A_\eta,B_\eta$ be matrices such that
$e^{tA_\eta}$, $e^{tB_\eta}$, and
$U_\eta(t)=e^{t(A_\eta+B_\eta)}$ are contractions for $t\geq0$.
This includes skew-Hermitian transport generators and Hermitian
negative semidefinite diffusion generators. Set
$S_\eta(s)=e^{sB_\eta}e^{sA_\eta}$.

\begin{assumption}[Uniform state-dependent commutator control]
For some $h_0>0$ and $K<\infty$, normalized initial vectors $u_{\eta,0}$
satisfy, on every grid,
\[
 \sup_{\substack{0\leq t\leq T\\0\leq r\leq s\leq h_0}}
 \norm{[B_\eta,A_\eta]e^{rB_\eta}e^{sA_\eta}
                  U_\eta(t)u_{\eta,0}}\leq K .
 \label{eq:hyp}
\]
\end{assumption}
This explicitly includes intermediate split states. Smoothness of the
continuum solution alone is not a proof of this uniform discrete
hypothesis.

\begin{proposition}[Uniform autonomous vector-norm splitting bound]
Under Assumption 1, for integer $L$ with $h=T/L\leq h_0$,
\[
 \norm{S_\eta(h)^L u_{\eta,0}-U_\eta(T)u_{\eta,0}}
 \leq \frac{KT^2}{2L}.
 \label{eq:theorem}
\]
If also $\norm{U_\eta(T)u_{\eta,0}}\geq m_*>0$ uniformly in $\eta$,
and the split output is nonzero, its normalized error is at most
$KT^2/(m_*L)$.
\end{proposition}
\begin{proof}
Suppress the grid subscript. Direct differentiation gives
\[
 S'(s)-(A+B)S(s)=[e^{sB},A]e^{sA}
 =\int_0^s e^{(s-r)B}[B,A]e^{rB}e^{sA}\,\dd r.
\]
Variation of constants provides the exact, remainder-free identity
\[
 (S(h)-U(h))v=
 \int_0^h\!\int_0^s
 U(h-s)e^{(s-r)B}[B,A]e^{rB}e^{sA}v\,\dd r\,\dd s .
\]
Use $v=U(jh)u_0$ and Assumption 1. Contractivity of the left factors
bounds each defect by $Kh^2/2$.
Finally the algebraic telescoping identity
\[
 S(h)^L-U(h)^L
 =\sum_{j=0}^{L-1}S(h)^{L-1-j}(S(h)-U(h))U(h)^j
\]
and contractivity bound the sum by $LKh^2/2=KT^2/(2L)$.
The normalized conclusion follows from Lemma 1.
\end{proof}

The assumption makes the conclusion conditional but precise: there is
no unspecified mesh-dependent Taylor remainder. For diffusion,
\eqref{eq:corrected-comm} indicates the necessary higher-derivative
quantities in a continuum analogue. Establishing uniform bounds on
the actual discrete partial trajectories remains a separate task.
For $d>2$, pairwise defects can be summed only after the corresponding
partial-product hypotheses are formulated. For time-dependent
generators, time-ordered propagators and coefficient-time regularity
must also be handled; neither extension is claimed here.

\section{Reproducible checks and interpretation}
The companion repository contains deterministic NumPy checks with
analytic references, invariant checks, and grid/step refinement.
They are classical computations; none executes a quantum circuit.
The transport example tests both conservation of $\norm f$ and
nonconservation of $\norm{\partial_yf}$.
The diffusion example checks the full commutator, the missing-term
ratio, and normalized split/exact behavior. Stencil checks verify
the extra zero mode and second-order consistency on a fixed smooth
Fourier mode. A negative-coefficient test illustrates the missing
positivity hypothesis.

\IfFileExists{results-table.tex}{\begin{table}[ht]
\centering\small
\begin{tabular}{rrrrrr}
\toprule
$T$ & $L$ & Normalized error & Error $L/T^2$ & Comparator & Ratio \\
\midrule
0.0002 & 64 & 1.1426e-07 & 182.819 & 38.062 & 4.803 \\
0.0002 & 128 & 5.7130e-08 & 182.815 & 38.062 & 4.803 \\
0.0002 & 256 & 2.8565e-08 & 182.813 & 38.062 & 4.803 \\
0.0002 & 512 & 1.4282e-08 & 182.812 & 38.062 & 4.803 \\
0.0005 & 64 & 5.8357e-07 & 149.394 & 38.351 & 3.895 \\
0.0005 & 128 & 2.9177e-07 & 149.388 & 38.351 & 3.895 \\
0.0005 & 256 & 1.4588e-07 & 149.385 & 38.351 & 3.895 \\
0.0005 & 512 & 7.2941e-08 & 149.383 & 38.351 & 3.895 \\
\bottomrule\end{tabular}
\caption{Classical normalized splitting check: $128\times128$ grid,
$p=1$, $a=0.2$, $b=0.5$, $m=4$, using the sine variant
$u_0=1+b\sin(2\pi m x)\sin(2\pi y)$. This has the same
analytic norm ratio as the cosine example. The comparator reconstructs
the printed prefactor with the initial continuum derivative norm
and the semidiscrete final norm. Finite numerical evidence supplements,
rather than proves, the analytic coefficient objection.}
\end{table}
}{
\emph{Run the companion checks for numerical values; the counterexamples
above are analytic and do not depend on this table.}}

The audit ledger records each finding's location, classification,
support, and proposed repair. Editorial issues include mixing integer
labels with dyadic coordinates on page 2, using $n$ for both total
and per-direction qubits, the time-order convention implicit in
Eq.~(6), and a primed/unprimed diffusion-prefactor inconsistency
on page 6. These issues should not be promoted to independent physical
or algorithmic refutations.

\section{Limitations and open work}
This note has not undergone external expert peer review and is
AI-assisted. Its strongest conclusions are the displayed identities,
analytic counterexamples, and the conditional matrix theorem with
a complete proof. The numerical checks supplement those arguments;
they do not prove a quantum advantage or an unconditional PDE theorem.
An exhaustive error claim would be unjustified.

A full replacement algorithm still needs a mesh-uniform consistency
and stability theorem for the chosen stencil, careful treatment of
normalized states and zero-mean data, a time-dependent splitting
analysis, explicit coefficient and state access models, non-unitary
success-cost bounds, observable definitions, and an end-to-end comparison
with appropriate classical algorithms. The counterexamples identify
repairs needed in the audited version without asserting that these
tasks cannot be completed.

\begin{thebibliography}{9}
\bibitem{original}
J. Zylberman, T. Fredon, N. F. Loureiro, and F. Debbasch,
\emph{Quantum algorithm for anisotropic diffusion and convection equations
with vector norm scaling}, arXiv:2603.08799v1 (2026).
\url{https://arxiv.org/abs/2603.08799v1}.
\bibitem{an}
D. An, D. Fang, and L. Lin,
\emph{Time-dependent unbounded Hamiltonian simulation with vector norm
scaling}, Quantum \textbf{5}, 459 (2021).
\url{https://doi.org/10.22331/q-2021-05-26-459}.
\bibitem{chapter}
J. Zylberman, T. Fredon, N. F. Loureiro, and F. Debbasch,
\emph{Quantum Algorithm for Anisotropic Diffusion and Convection Equations
with Vector Norm Scaling}, conference chapter, pp.~255--263.
Public source-linked version of record; full text not audited here.
\url{https://doi.org/10.1007/978-3-032-13855-2_23}.
\end{thebibliography}
\end{document}
